\documentclass{article}
\input{../../lib.sty}

\title{\texttt{beta\_via\_approx\_dp}}
\author{\AishwaryaRamasethuAuthor, \YuJuKuAuthor, \JordanAwanAuthor, \MichaelShoemateAuthor}
\begin{document}

\maketitle

\contrib

\section{Hoare Triple}

\subsection*{Precondition}
\subsubsection*{Compiler-verified}
\begin{itemize}
    \item Arguments \texttt{epsilon}, \texttt{delta}, and \texttt{alpha} of type \texttt{f64}
\end{itemize}

\subsubsection*{Caller-verified}
The source parameters satisfy finite nonnegative $\epsilon$, $\delta \in [0,1]$,
and $\alpha \in [0,1]$. The function must complete successfully.

\subsection*{Pseudocode}

\lstinputlisting[language=Python,firstline=2,escapechar=`]{./pseudocode/approxdp_beta.py}

\subsection*{Postcondition}

\begin{theorem}
    For all $\alpha \in [0, 1]$, \texttt{beta\_via\_approx\_dp} returns a value
    no greater than the exact approximate-DP tradeoff curve
    $f_{\epsilon,\delta}(\alpha)$.
\end{theorem}

\begin{definition}
    \label{def-approxdp-point-beta}
    For a fixed approximate-DP certificate with parameters $\epsilon$ and $\delta$, define
    \begin{equation}
        f_{\epsilon,\delta}(\alpha) = \max(0, 1 - \delta - \exp(\epsilon) \alpha,
        \exp(-\epsilon) \max(1 - \delta - \alpha, 0)).
    \end{equation}
\end{definition}

\begin{proof}
The function directly computes exact rational conversions of directed interval
endpoints satisfying
\[
    \texttt{exp\_eps\_up} \ge \exp(\epsilon), \qquad
    \texttt{exp\_neg\_eps\_down} \le \exp(-\epsilon).
\]
The pseudocode uses these coefficients with exact rational arithmetic.

For $\alpha \in [0,1]$, the first affine branch is no greater than
$1 - \delta - \exp(\epsilon)\alpha$ because \texttt{exp\_eps\_up} is an
upper bound and $\alpha$ is nonnegative. The second branch is similarly no
greater than its exact counterpart because \texttt{exp\_neg\_eps\_down}
is a lower bound. Taking the maximum with zero preserves the pointwise lower
bound. The final conversion to \texttt{f64} rounds downward. The endpoints return
zero at $\alpha=1$ and a downward-rounded $1-\delta$ at $\alpha=0$,
which satisfy the same bound.
\end{proof}

\end{document}
